% I. Anexo técnico: software (Guía apdo. 3.I)
% Fuentes de verdad: docs/DATA_SCHEMA.md, device/device.ino, server/mqtt_logger.py
\chapter{Anexo técnico: software y formato de datos}
\label{anexo:software}

\section{Formato de la trama de telemetría}

El nodo publica en el tema \texttt{fridge/sensors} un objeto JSON con nueve variables por
ciclo de medida. El listado~\ref{listado:trama} muestra un ejemplo de trama.

\begin{lstlisting}[caption=Ejemplo de trama de telemetria publicada por el nodo.,label=listado:trama]
{"tempExt":4.31,"accX":0.12,"accY":-0.05,"accZ":9.79,
 "gyroX":0.01,"gyroY":0.00,"gyroZ":-0.02,"motorTemp":38.6,"noise":142}
\end{lstlisting}

El cuadro de las nueve variables de este canal, con su unidad e intervalo, figura en el
capítulo~\ref{cap:diseno}. El canal de ráfaga, publicado en el tema
\texttt{fridge/vibration}, transporta en cambio características ya calculadas y no medidas
instantáneas. El cuadro~\ref{table:rafaga} recoge su contenido completo.

\begin{table}[ht!]
\caption{Campos del canal de ráfaga. El sufijo \texttt{\_x}, \texttt{\_y} o \texttt{\_z}
identifica el eje.}
\label{table:rafaga}
\centering
\begin{tabular}{|p{2.6cm}|p{1.5cm}|p{8.3cm}|}
\hline
\textbf{Campo} & \textbf{Unidad} & \textbf{Contenido} \\
\hline
\texttt{vib\_fs} & Hz & Frecuencia de muestreo de la ráfaga (1000) \\
\hline
\texttt{vib\_n} & --- & Muestras por ráfaga (1024) \\
\hline
\texttt{ms\_capture} & ms & Duración real de la captura. Es el control de calidad de la
ráfaga: una desviación respecto de \mbox{1024 ms} comprime el eje de frecuencias \\
\hline
\texttt{ms\_total} & ms & Captura más cálculo de características \\
\hline
\texttt{failed\_bursts} & --- & Ráfagas descartadas acumuladas \\
\hline
\texttt{bad\_frames} & --- & Tramas del canal lento descartadas acumuladas \\
\hline
\texttt{retries} & --- & Reintentos de lectura en la ráfaga en curso \\
\hline
\texttt{total\_retries} & --- & Reintentos acumulados desde el arranque \\
\hline
\texttt{rms\_*} & m/s$^2$ & Valor eficaz de la componente alterna \\
\hline
\texttt{peak\_*} & m/s$^2$ & Máximo absoluto tras eliminar la continua \\
\hline
\texttt{kurt\_*} & --- & Kurtosis. 3 para ruido gaussiano, 1,5 para senoide pura \\
\hline
\texttt{fdom\_*} & Hz & Frecuencia dominante, refinada por interpolación parabólica \\
\hline
\texttt{adom\_*} & m/s$^2$ & Amplitud estimada en la frecuencia dominante \\
\hline
\texttt{aud\_fs} & Hz & Frecuencia de muestreo del audio (16000) \\
\hline
\texttt{aud\_n} & --- & Muestras de audio analizadas \\
\hline
\texttt{aud\_rms} & --- & Valor eficaz de la señal acústica, en unidades relativas \\
\hline
\texttt{aud\_b0} a \texttt{aud\_b3} & --- & Fracción de energía espectral en las bandas
0--250, 250--1000, 1000--4000 y 4000--8000 \mbox{Hz}. Las cuatro suman la unidad \\
\hline
\end{tabular}
\end{table}

Las cuatro bandas acústicas expresan fracciones y no niveles absolutos: describen cómo se
reparte la energía del sonido, no cuánta hay. El nivel absoluto reside en \texttt{aud\_rms}.

\section{Valores centinela y tratamiento de datos ausentes}

El sistema distingue explícitamente entre una medida válida y la ausencia de medida. Un
valor de \mbox{$-127$ \textdegree C} en la temperatura externa indica sonda desconectada, y
un bloque de siete variables del acelerómetro con valor exactamente nulo indica que el bus
I\textsuperscript{2}C no pudo recuperarse en ese ciclo. Ambos casos deben filtrarse antes
de calcular cualquier estadístico, ya que su inclusión distorsiona los promedios del
conjunto de datos.

\section{Dependencias del entorno de desarrollo}

El cuadro~\ref{table:dependencias} recoge las versiones empleadas, tomadas del registro de
compilación del firmware y del fichero de dependencias del concentrador. Se documentan de
forma explícita porque el comportamiento de una de ellas resultó determinante para el diseño
de la adquisición, según se describe en el capítulo~\ref{cap:diseno}.

\begin{table}[ht!]
\caption{Versiones del entorno de desarrollo y de las dependencias.}
\label{table:dependencias}
\centering
\begin{tabular}{|p{5.4cm}|p{2.4cm}|p{4.4cm}|}
\hline
\textbf{Componente} & \textbf{Versión} & \textbf{Función} \\
\hline
Núcleo Arduino para ESP32 & 3.3.1 & Entorno de ejecución del nodo \\
\hline
Wire & 3.3.0 & Bus I\textsuperscript{2}C \\
\hline
Adafruit MPU6050 & 2.2.9 & Configuración del acelerómetro \\
\hline
Adafruit BusIO & 1.17.4 & Abstracción de bus \\
\hline
Adafruit Unified Sensor & 1.1.15 & Interfaz común de sensores \\
\hline
OneWire & 2.3.8 & Bus 1-Wire \\
\hline
DallasTemperature & 4.0.6 & Sonda DS18B20 \\
\hline
WiFi / Networking & 3.3.0 & Conectividad del nodo \\
\hline
PubSubClient & 2.8 & Cliente MQTT del nodo \\
\hline
Python & 3.11 o superior & Entorno del concentrador \\
\hline
paho-mqtt & 2.1.0 & Cliente MQTT del registrador \\
\hline
python-dotenv & 1.0.1 & Configuración por variables de entorno \\
\hline
Mosquitto & 2.0 o superior & Intermediario de mensajería \\
\hline
\end{tabular}
\end{table}

La biblioteca Adafruit MPU6050 se emplea únicamente para la configuración inicial del
sensor, y no para la lectura de medidas. El motivo es que su función de obtención de datos
devuelve un indicador de éxito incondicional y opera sobre un almacén no inicializado, de
modo que ante un fallo del bus entrega como medida el contenido residual de la pila. La
lectura se realiza por ello mediante acceso directo a los registros del dispositivo, con la
validación descrita en el capítulo~\ref{cap:diseno}. Se documenta esta circunstancia porque
condiciona la reproducibilidad: una implementación que confiara en esa biblioteca para la
adquisición reproduciría el mismo defecto.

\section{Aprovisionamiento del concentrador}

El capítulo~\ref{cap:metodologia} expone los criterios que rigen el despliegue; esta sección
recoge el procedimiento reproducible. El aprovisionamiento está automatizado en un
\textit{script} que admite ejecutar las operaciones de una vez o por fases separadas, según
el listado~\ref{lst:aprovisionamiento}.

\begin{lstlisting}[caption={Aprovisionamiento completo del concentrador.}, label={lst:aprovisionamiento}]
# Desde el equipo de desarrollo: copiar el proyecto al concentrador.
# Se excluyen el entorno virtual, la configuracion local y los datos
# capturados, para no sobrescribir lo que ya reside en el concentrador.
rsync -av --exclude .venv --exclude 'server/.env' --exclude 'server/data' \
      <ruta-local>/ <usuario>@<direccion>:/home/<usuario>/TFM/

# En el concentrador: aprovisionamiento completo
cd ~/TFM
sudo AP_PASS='<clave-de-la-red>' ./server/provision-pi.sh todo
\end{lstlisting}

El cuadro~\ref{table:fases-aprovisionamiento} detalla las operaciones y su orden, que no es
arbitrario.

\begin{table}[ht!]
\caption{Fases del aprovisionamiento del concentrador.}
\label{table:fases-aprovisionamiento}
\centering
\begin{tabular}{|p{2.1cm}|p{10.4cm}|}
\hline
\textbf{Fase} & \textbf{Operación} \\
\hline
\texttt{paquetes} & Instalación del intermediario de mensajería y del entorno de ejecución
de Python. Requiere acceso a Internet \\
\hline
\texttt{broker} & Configuración del intermediario, con examen previo de la configuración
existente \\
\hline
\texttt{logger} & Entorno virtual, configuración por variables de entorno y declaración del
registrador como servicio del sistema. Incluye una comprobación de extremo a extremo con una
trama sintética \\
\hline
\texttt{ap} & Conversión de la interfaz inalámbrica en punto de acceso. Última operación \\
\hline
\texttt{comprobar} & Verificación de la cadena completa. No requiere privilegios \\
\hline
\texttt{revertir} & Restitución de la interfaz a modo cliente y detención del registrador,
conservando los datos capturados \\
\hline
\end{tabular}
\end{table}

Tres particularidades del procedimiento responden a restricciones del entorno. La primera es
que las comprobaciones previas se completan antes de modificar el sistema, de modo que una
condición no satisfecha aborta la ejecución sin dejar el equipo en un estado intermedio. La
segunda es que la activación del punto de acceso interrumpe la red por la que llega la sesión
de administración remota, por lo que constituye la última operación y se lanza desacoplada de
la sesión mediante el gestor de servicios, con objeto de que se complete aunque la conexión se
pierda; los datos de acceso se escriben en un fichero del concentrador antes de esa
activación, para que puedan consultarse localmente después. La tercera es la existencia de la
fase de reversión, exigible cuando el equipo presta además otros servicios: la
restitución no elimina los ficheros de datos capturados.

La instalación exige acceso a Internet, mientras que el punto de acceso lo suprime en la
interfaz inalámbrica. Ese es el motivo del orden establecido, y no una preferencia
organizativa: invertirlo dejaría el concentrador sin posibilidad de completar la instalación
de dependencias.

\section{Puesta en marcha y verificación del concentrador}

El concentrador no requiere intervención para entrar en servicio. El intermediario de
mensajería, el registrador y el punto de acceso inalámbrico están declarados como servicios
del sistema con arranque automático, de modo que la alimentación eléctrica es la única
condición necesaria. Esta propiedad se estableció como requisito porque una campaña de
adquisición prolongada debe sobrevivir a un corte de suministro sin perder más datos que los
del intervalo sin corriente, y el registrador escribe en modo de adición, con lo que
prosigue el fichero de la jornada en lugar de sustituirlo.

El registrador se declara además con reinicio automático ante terminación anómala. El
criterio subyacente es que los datos en bruto son irrecuperables si se pierden, de modo que
un fallo puntual del proceso no debe interrumpir la campaña.

La verificación de la cadena completa se realiza con el procedimiento recogido en el
listado~\ref{lst:verificacion}, que recorre el camino del dato en el mismo orden en que este
fluye. Esa ordenación es deliberada: el primer eslabón que falla localiza el corte, de forma
que no es necesario examinar los posteriores.

\begin{lstlisting}[caption={Verificacion de la cadena de adquisicion en el concentrador.}, label={lst:verificacion}]
# Verificacion completa: punto de acceso, nodo asociado, intermediario,
# tramas de ambos canales, registrador y crecimiento de los ficheros
./server/provision-pi.sh comprobar

# Seguimiento durante la campana
journalctl -u tfm-logger -f
wc -l server/data/*.csv
\end{lstlisting}

La comprobación del registrador no se limita a constatar la existencia de los ficheros de
datos, sino que contrasta su número de filas en dos instantes separados. La distinción es
pertinente: un fichero presente que no crece constituye un modo de fallo silencioso, y se
observó durante la puesta a punto que puede producirse con el servicio aparentemente activo y
las tramas llegando correctamente al intermediario.

El origen de aquel incidente ilustra un aspecto del comportamiento del sistema operativo que
conviene documentar. Al eliminar los ficheros de datos con el registrador en ejecución, el
descriptor abierto por el proceso continúa referenciando el fichero original, ya desvinculado
del sistema de ficheros, por lo que la escritura prosigue sobre un destino inaccesible
mientras el fichero visible permanece vacío. La reapertura solo se produce al cambiar la
fecha. Se establece por ello como procedimiento que toda manipulación del directorio de datos
se realice con el servicio detenido, según el listado~\ref{lst:limpieza}.

\begin{lstlisting}[caption={Reinicio del directorio de datos antes de una campana.}, label={lst:limpieza}]
D=server/data/descartes-$(date +%Y%m%d-%H%M)
sudo systemctl stop tfm-logger        # detener ANTES de tocar los ficheros
mkdir -p "$D" && mv server/data/*.csv "$D"/
sudo systemctl start tfm-logger       # los recrea con su cabecera
\end{lstlisting}

Los ficheros se archivan en lugar de eliminarse, dado que las capturas de puesta a punto
constituyen la evidencia de las incidencias descritas en el capítulo~\ref{cap:diseno}. Debe
evitarse crear los ficheros de forma manual tras retirarlos: el registrador escribe la
cabecera únicamente cuando el fichero no existe o está vacío, y un fichero preexistente
produciría un conjunto de datos sin nombres de columna.

\section{Procedimiento de despliegue del modelo en el nodo}

El modelo ajustado llega al firmware a través de un fichero de cabecera generado por el entorno
de análisis, según la secuencia siguiente:

\begin{enumerate}
\item El procedimiento de exportación ajusta el modelo sobre el conjunto de referencia y escribe
      \texttt{device/modelo\_referencia.h}, que contiene los parámetros y, en su encabezado, la
      campaña, la duración, el número de ráfagas y episodios, los umbrales del filtro de calidad
      y la fecha de generación.
\item Un segundo procedimiento exporta las ráfagas reales con la puntuación que produce la
      biblioteca de referencia sobre cada una, para contrastar la implementación embarcada.
\item Las rutinas del detector se compilan y verifican en un computador frente a esos valores.
\item Solo entonces se compila el firmware y se programa la placa.
\end{enumerate}

La separación entre los parámetros y las rutinas es deliberada: reentrenar el modelo consiste en
volver a ejecutar la exportación y recompilar, sin modificar código. El fichero generado no se
edita a mano, y su encabezado permite rastrear cualquier firmware programado en una placa hasta
los datos que lo produjeron.

\section{Especificidad del modelo respecto del activo}

Procede consignar una restricción de despliegue que el propio procedimiento de verificación puso
de manifiesto. El modelo exportado incorpora las medianas del activo sobre el que se ajustó, su
umbral de discriminación entre marcha y parada y su referencia de estado nominal. Programarlo en
un nodo instalado sobre otro activo produce la declaración de medida no evaluable en la
totalidad de las ráfagas: se comprobó que el activo con fallo presenta un valor eficaz en marcha
de \mbox{0,16 m/s$^2$} frente al umbral de \mbox{0,199 m/s$^2$} derivado del activo de
referencia, de modo que sus ráfagas se clasifican como capturadas con el compresor detenido.

\textbf{El modelo es específico del activo}, y cada nodo requiere su propia campaña de
referencia. De ello se sigue asimismo que el nodo necesita una fase de referencia antes de poder
emitir su primer veredicto, dado que las características normalizadas por el propio activo
exigen conocer su mediana.

\section{Rutina de recuperación del bus}

El listado~\ref{lst:recuperacion} recoge la rutina completa. Se invoca cuando una lectura ha
fallado de forma efectiva, y no de forma preventiva en cada ciclo.

\begin{lstlisting}[language=C++, caption={Rutina de recuperacion del bus I2C.}, label={lst:recuperacion}]
bool recoverI2C() {
  Wire.end();
  delay(100);
  Wire.begin(I2C_SDA, I2C_SCL);
  Wire.setClock(I2C_FREQ_HZ);

  if (mpu.begin()) {
    mpu.setAccelerometerRange(MPU6050_RANGE_4_G);
    mpu.setFilterBandwidth(MPU6050_BAND_260_HZ);
    mpu.setSampleRateDivisor(0);
    delay(50);            // los registros valen 0 hasta la 1a conversion
    updateGyroScale();    // begin() reprograma el rango del giroscopo
    return true;
  }
  return false;
}
\end{lstlisting}

Tres aspectos del listado responden a incidencias detectadas en la placa. Los pines del bus
se indican de forma explícita en la reapertura, en lugar de heredarse de la configuración de
la variante de la placa, para que la recuperación no reasigne el bus a pines distintos de los
cableados. La espera posterior a la reinicialización es necesaria porque la puesta en marcha
del sensor conlleva un reinicio del dispositivo, tras el cual sus registros de medida
permanecen a cero hasta que concluye la primera conversión; leerlos antes introduciría
muestras nulas indistinguibles de una medida. Por último, la escala del giroscopio se vuelve
a consultar al dispositivo, dado que la reinicialización reprograma su fondo de escala y
suponerlo en lugar de leerlo alteraría de forma silenciosa las unidades de tres variables del
conjunto de datos.
